\documentclass[11pt,a4paper]{article}

% =====================================================
% preamble.tex — Configuración común de estilo
% Para enunciados de ejercitación en Tekla Structures
% =====================================================

% --- Codificación e idioma ---
\usepackage[utf8]{inputenc}
\usepackage[T1]{fontenc}
\usepackage[spanish,es-tabla,es-noshorthands]{babel}
\usepackage{lmodern}
\usepackage{microtype}

% --- Geometría de página ---
\usepackage[a4paper,margin=2.5cm,headheight=14pt]{geometry}

% --- Matemática y unidades ---
\usepackage{amsmath,amssymb}
\usepackage{siunitx}
\sisetup{
  output-decimal-marker={,},
  per-mode=symbol,
  group-separator={.},
  group-minimum-digits=4
}

% --- Tablas ---
\usepackage{booktabs}
\usepackage{array}
\usepackage{multirow}
\usepackage{tabularx}
\usepackage{longtable}

% --- Figuras ---
\usepackage{graphicx}
\usepackage{caption}
\usepackage{subcaption}
\usepackage{float}

% --- Gráficos vectoriales ---
\usepackage{tikz}
\usetikzlibrary{arrows.meta,positioning,calc,patterns,shapes.geometric}

% --- Listas ---
\usepackage{enumitem}
\setlist{itemsep=2pt,parsep=2pt}

% --- Color y cajas ---
\usepackage{xcolor}
\definecolor{teklablue}{RGB}{0,87,156}
\definecolor{boxgray}{RGB}{245,245,245}
\definecolor{accentorange}{RGB}{230,120,30}

\usepackage[most]{tcolorbox}

% Caja para comandos / atajos
\newtcolorbox{comando}{
  colback=boxgray,
  colframe=teklablue,
  boxrule=0.5pt,
  arc=2pt,
  left=6pt, right=6pt, top=4pt, bottom=4pt,
  fonttitle=\bfseries,
  title=Comando Tekla
}

% Caja para tips / consejos
\newtcolorbox{tip}{
  colback=yellow!8,
  colframe=accentorange,
  boxrule=0.5pt,
  arc=2pt,
  fonttitle=\bfseries,
  title=Tip
}

% Caja para entregable
\newtcolorbox{entregable}{
  colback=green!5,
  colframe=green!50!black,
  boxrule=0.5pt,
  arc=2pt,
  fonttitle=\bfseries,
  title=Entregable
}

% Caja para advertencia
\newtcolorbox{advertencia}{
  colback=red!5,
  colframe=red!70!black,
  boxrule=0.5pt,
  arc=2pt,
  fonttitle=\bfseries,
  title=Atención
}

% --- Código y comandos en línea ---
\usepackage{listings}
\lstset{
  basicstyle=\ttfamily\small,
  breaklines=true,
  frame=single,
  backgroundcolor=\color{boxgray},
  rulecolor=\color{teklablue}
}

% Comando para resaltar comandos/menús de Tekla en línea
% Uso: \tk{File > Open}  o  \tk{Component Catalog}
\newcommand{\tk}[1]{\textsf{\textbf{\textcolor{teklablue}{#1}}}}

% Comando para nombres de archivo / rutas
\newcommand{\arch}[1]{\texttt{#1}}

% Comando para teclas / shortcuts
% Uso: \key{Ctrl+S}
\newcommand{\key}[1]{\fbox{\footnotesize\texttt{#1}}}

% --- Encabezado y pie de página ---
\usepackage{fancyhdr}
\pagestyle{fancy}
\fancyhf{}
\fancyhead[L]{\small\textit{\tituloEjercicio}}
\fancyhead[R]{\small\thepage}
\fancyfoot[C]{\small Tekla Structures 2022 — Nivel avanzado}
\renewcommand{\headrulewidth}{0.4pt}
\renewcommand{\footrulewidth}{0.4pt}

% --- Hipervínculos y referencias ---
\usepackage{hyperref}
\hypersetup{
  colorlinks=true,
  linkcolor=teklablue,
  urlcolor=teklablue,
  citecolor=teklablue,
  pdfborder={0 0 0}
}
\usepackage[noabbrev,spanish]{cleveref}

% --- Formato de secciones ---
\usepackage{titlesec}
\titleformat{\section}
  {\Large\bfseries\color{teklablue}}
  {\thesection.}{0.5em}{}
\titleformat{\subsection}
  {\large\bfseries\color{teklablue!80!black}}
  {\thesubsection.}{0.5em}{}

% =====================================================
% commands.tex — Variables propias de cada enunciado
% Editar estos campos al inicio para cada ejercicio
% =====================================================

% --- Identificación del ejercicio ---
\newcommand{\tituloEjercicio}{Ejercicio: [TÍTULO]}
\newcommand{\subtituloEjercicio}{[Subtítulo descriptivo]}
\newcommand{\codigoEjercicio}{TEKLA-XXX}
\newcommand{\nivelEjercicio}{Avanzado}
\newcommand{\duracionEjercicio}{[X] horas}

% --- Datos del curso ---
\newcommand{\nombreCurso}{[Nombre del curso]}
\newcommand{\autorEnunciado}{[Docente]}
\newcommand{\fechaEjercicio}{\today}

% --- Versión del software ---
\newcommand{\versionTekla}{Tekla Structures 2022}
\newcommand{\entornoTekla}{[Environment / Configuration]}

% --- Atajos para terminología recurrente ---
\newcommand{\CC}{\tk{Component Catalog}}
\newcommand{\AC}{\tk{Applications \& Components}}
\newcommand{\OB}{\tk{Object Browser}}
\newcommand{\OR}{\tk{Object Representation}}
\newcommand{\OG}{\tk{Object Group}}
\newcommand{\SF}{\tk{Selection Filter}}
\newcommand{\VF}{\tk{View Filter}}


% =====================================================
% Editar en commands.tex:
%   \tituloEjercicio, \subtituloEjercicio, etc.
% =====================================================

\renewcommand{\tituloEjercicio}{Ejercicio de Hormigón Armado}
\renewcommand{\subtituloEjercicio}{[Subtítulo: tipo de estructura]}
\renewcommand{\codigoEjercicio}{TEKLA-HA-01}

\begin{document}

% =====================================================
% PORTADA
% =====================================================
\begin{titlepage}
\centering
\vspace*{2cm}
{\Huge\bfseries\color{teklablue} \tituloEjercicio \par}
\vspace{0.5cm}
{\Large \subtituloEjercicio \par}
\vspace{2cm}
\rule{0.6\textwidth}{0.4pt}\par
\vspace{1cm}
\begin{tabular}{rl}
\textbf{Código:} & \codigoEjercicio \\
\textbf{Curso:} & \nombreCurso \\
\textbf{Nivel:} & \nivelEjercicio \\
\textbf{Duración estimada:} & \duracionEjercicio \\
\textbf{Software:} & \versionTekla \\
\textbf{Docente:} & \autorEnunciado \\
\textbf{Fecha:} & \fechaEjercicio \\
\end{tabular}
\vfill
\end{titlepage}

\tableofcontents
\newpage

% =====================================================
\section{Objetivos de aprendizaje}
% =====================================================
Al finalizar este ejercicio, el alumno será capaz de:
\begin{itemize}
    \item [Objetivo 1: ej. modelar elementos in-situ con \tk{Concrete Beam} y \tk{Concrete Slab}]
    \item [Objetivo 2: ej. aplicar \tk{Rebar Sets} con distribuciones complejas]
    \item [Objetivo 3: ej. generar \tk{Pour Units} y \tk{Pour Breaks}]
    \item [Objetivo 4: ej. crear un \tk{Custom Component} paramétrico]
    \item [Objetivo 5: ej. producir planos de armado y planillas de hierros]
\end{itemize}

% =====================================================
\section{Descripción del proyecto}
% =====================================================
[Descripción breve, 1--2 párrafos, de qué se va a modelar y en qué contexto profesional se ubica.]

% =====================================================
\section{Datos de partida}
% =====================================================
\subsection{Geometría general}
\begin{itemize}
    \item [Dimensiones generales]
    \item [Cantidad de niveles / vanos]
    \item [Alturas y luces]
\end{itemize}

\subsection{Materiales}
\begin{table}[H]
\centering
\begin{tabular}{lll}
\toprule
\textbf{Elemento} & \textbf{Hormigón} & \textbf{Acero} \\
\midrule
[Vigas]    & [H-XX] & [ADN 420] \\
[Columnas] & [H-XX] & [ADN 420] \\
[Losas]    & [H-XX] & [ADN 420] \\
\bottomrule
\end{tabular}
\caption{Materiales del proyecto.}
\end{table}

\subsection{Normativa de referencia}
\begin{itemize}
    \item [CIRSOC 201 / ACI 318 / norma local]
    \item [Estándar de representación gráfica]
\end{itemize}

% =====================================================
\section{Esquema de referencia}
% =====================================================
[Insertar planta, cortes y vistas. Espacio para imagen o TikZ.]

\begin{figure}[H]
\centering
\fbox{\parbox[c][6cm][c]{0.8\textwidth}{\centering\textit{[Esquema de referencia: planta y cortes]}}}
\caption{Esquema general del proyecto.}
\label{fig:esquema}
\end{figure}

% =====================================================
\section{Consigna paso a paso}
% =====================================================

\subsection{Configuración inicial del modelo}
\begin{enumerate}
    \item Crear un nuevo modelo con el entorno \entornoTekla.
    \item Configurar \tk{Project Properties} con [datos del proyecto].
    \item Definir \tk{Grids} y \tk{Levels} según el esquema de referencia.
\end{enumerate}

\begin{tip}
[Tip sobre nomenclatura de grids o configuración inicial]
\end{tip}

\subsection{Modelado de elementos primarios}
\begin{enumerate}
    \item [Modelar columnas con \tk{Concrete Column}]
    \item [Modelar vigas con \tk{Concrete Beam}]
    \item [Modelar losas con \tk{Concrete Slab}]
\end{enumerate}

\subsection{Modelado de elementos secundarios}
\begin{enumerate}
    \item [Muros, escaleras, fundaciones, etc.]
\end{enumerate}

\subsection{Armaduras}
\begin{enumerate}
    \item [Aplicar \tk{Rebar Set} en vigas]
    \item [Aplicar \tk{Rebar Set} en columnas]
    \item [Distribuciones especiales en zonas de empalme]
\end{enumerate}

\begin{comando}
[Sintaxis o flujo específico de un comando clave]
\end{comando}

\subsection{Custom Component}
\begin{enumerate}
    \item Crear un \tk{Custom Component} para [detalle específico].
    \item Parametrizar con [variables].
    \item Validar con un caso de uso.
\end{enumerate}

\subsection{Pours y numeración}
\begin{enumerate}
    \item Definir \tk{Pour Breaks} según [criterio].
    \item Ejecutar \tk{Numbering > Numbering Series}.
    \item Verificar marcas con \tk{Inquire}.
\end{enumerate}

\subsection{Planos y reportes}
\begin{enumerate}
    \item Generar \tk{General Arrangement Drawings} de cada nivel.
    \item Generar \tk{Cast Unit Drawings} para [elementos].
    \item Producir \tk{Reports} con planilla de hierros.
\end{enumerate}

% =====================================================
\section{Entregables}
% =====================================================
\begin{entregable}
\begin{enumerate}
    \item Archivo de modelo \arch{.db1} comprimido en \arch{.zip}.
    \item Planos de encofrado en PDF (un PDF por nivel).
    \item Planos de armado en PDF.
    \item Planilla de hierros en formato Excel o CSV.
    \item Documento con capturas que muestren [hitos clave del modelado].
\end{enumerate}
\end{entregable}

% =====================================================
\section{Criterios de evaluación}
% =====================================================
\begin{table}[H]
\centering
\begin{tabularx}{\textwidth}{Xcc}
\toprule
\textbf{Criterio} & \textbf{Puntaje} & \textbf{Obtenido} \\
\midrule
Geometría correcta del modelo & 20 & \\
Aplicación adecuada de rebar sets & 25 & \\
Calidad del custom component & 20 & \\
Numeración y pours coherentes & 10 & \\
Calidad de planos generados & 15 & \\
Completitud de entregables & 10 & \\
\midrule
\textbf{Total} & \textbf{100} & \\
\bottomrule
\end{tabularx}
\end{table}

% =====================================================
\section{Tips y errores comunes}
% =====================================================
\begin{tip}
[Tip 1]
\end{tip}

\begin{advertencia}
[Error frecuente y cómo evitarlo]
\end{advertencia}

% =====================================================
\section{Referencias y documentación}
% =====================================================
\begin{itemize}
    \item Tekla User Assistance: \url{https://support.tekla.com/}
    \item [Otras referencias]
\end{itemize}

\end{document}
